% a06_BALtorus_d5: part of Step 3 of the proof of thm:a06:T1 (input from sec/a06_BALtorus_c.tex).
% Group (C), second half: the Z-factor expansions of Ba2 p. 578 (the terms containing B~^(k)) and the eps_0-freedom of the per-step large-field constants.

\subsubsection*{Step 3(C$_2$): the $Z^{(j)}$ expansions of \cite[p.~578]{Ba2}}

At step $k+1$ the fluctuation determinants $Z^{(j)}$, $j<k$, of \cite[(2.92)]{Ba2} are expanded in the new field. With $\tilde B$ the regular part of the configuration \cite[(2.93)--(2.94)]{Ba2}, $A_1\defeq\theta_{k+1}A'^{(k)}$ and $A_2\defeq\tilde B^{(k)}$, both rescaled to scale $j$, and $F_j(Y)\defeq\ln Z^{(j)}(Y)=-\frac12\ln\det X^{(j)}(Y)+\mathrm{const}$, $X^{(j)}(Y)\defeq\Lambda(\Delta^{(j)}(\Omega,Y)+aL^{-2}P(Y))\Lambda$ on $\ell^2(\Lambda;\R^N)$, $\Omega\defeq B^j(\Lambda^{(j)}_2)$, $\Lambda\defeq\Lambda^{(j)}_5$ \cite[(2.92), (2.99)]{Ba2}, exactly
\[
F_j(\tilde B+A_1+A_2)=F_j(\tilde B)+P_{\tilde B}[A_1]+R_j+D_j,\qquad R_j\defeq F_j(\tilde B+A_1)-F_j(\tilde B)-P_{\tilde B}[A_1],\quad D_j\defeq F_j(\tilde B+A_1+A_2)-F_j(\tilde B+A_1),
\]
$P_{\tilde B}[A_1]$ the order-$\bar n$ Taylor polynomial \cite[(2.104)]{Ba2}. $R_j$ (the order-$(\bar n+1)$ trace term) is bounded by the Cauchy estimate on the analytic family of \cref{lem:a06d:S1} below: $\sum_{j<k}|R_j|\le C_{\kappa}t_k^\kappa|\Lambda^{(k)}_2|/(L^{\bar n+1-d}-1)$ for $\kappa<(\bar n+1)/2$. The terms $D_j$ ``containing $\tilde B^{(k)}$'', which \cite[p.~578]{Ba2} bounds by $O((L^j\eps)^\kappa)|B^{k-j}(\Lambda^{(k)}_2)|$ without argument, are the content of this part: a per-site count at scale $j$ gives only $Ct_k^{\kappa+1/2}L^{(k-j)(d-1)}|\Lambda^{(k)}_2|$, which is not summable over $j<k$. The mechanism is that $A_2$ is smooth at scale $k$, so at scale $j$ it is almost pure gauge on every region of size $O(1)$ (its curvature is $L^{-2(k-j)}$ times smaller than at scale $k$, its size only $L^{-(k-j)}$ times smaller), and so is $B\defeq\tilde B+A_1$; $Z^{(j)}$ is exactly gauge invariant and even under charge conjugation, so after a local gauge fixing the first variation of each local piece of $\ln Z^{(j)}$ is bilinear in the two curvatures, $L^{-2(k-j)}\times L^{-2(k-j)}$ per scale-$j$ unit cube against $L^{3(k-j)}$ cubes per scale-$k$ cube, leaving $L^{-(k-j)}$, which sums. No identity on a single graph slot is used: the local pieces are exact gauge-invariant functions.

\paragraph{Setting.} Fix a scale $j$, $\Omega$ a union of unit blocks of the fine lattice $\xi\Z^3$, $\xi=L^{-j}$ (a box or a torus, Neumann bonds), $\Lambda$ a union of $L$-blocks in $\Omega^{(j)}$, $e_j\defeq et_j^{1/2}$, $U(z)\defeq\exp(qe_j\xi z)$ for $z\in\C$ ($q$ real antisymmetric, so $U$ is entire, $U(z)^T=U(-z)=U(z)^{-1}$, orthogonal for real $z$). For a real bond field $B$ on $\Omega$ put $X(B)\defeq X_\Lambda(\Omega,B)$ as above and $F(B)\defeq-\frac12\ln\det X(B)$. For a complex bond field $Z$ let $\hat w_\Delta(Z)\defeq\norm qe_j\sup\{|Z_b|:b_-\in\Delta\}$ for a unit cube $\Delta$ and $w(Z)\defeq\sup_\Delta\hat w_\Delta(Z)$; for a real $B$ let $f(B)\defeq\norm qe_j\sup|dB|$ be its curvature amplitude. Hypotheses: (Cur) $f(B)\le c_{\mathrm{cur}}$, with $c_{\mathrm{cur}}>0$ (depending on $d,\norm q,a,L$) such that the defects $2c_p\eps_g+d\kappa^2$ of the cube coercivity (\cref{lem:a06A:cube}) are at most half of $\min\{2\ell^{-2},c_p\}$ for $\ell\in\{1,L\}$; \cref{lem:a06A:cube} uses the field only through $|F|$, so under (Cur) it and \cref{lem:a06A:coer,lem:a06A:CT,lem:a06A:blockinv} and the first bound of \cref{prop:a06A:Ba1P22} hold for $B$ with no other condition, in particular $\gamma_X\le X(B)\le\gamma_1\defeq a+aL^{-2}$ and $|\Delta^{(j)}(\Omega,B;y,y')|\le c_7e^{-\delta_7|y-y'|}$ \cite[Prop.~2.2, (2.27)]{Ba1}, hence $|X(B)(y,y')|\le c_Xe^{-\delta_X|y-y'|}$ with $c_X\defeq c_7+aL^{-2}e^{\delta_7L}$, $\delta_X\defeq\delta_7$. All constants below depend only on $(d,N,\norm q,a,L,M)$ and the constants of \cref{thm:a06A:groupA}, never on $j,k,\eps,\eps_0,\Omega,\Lambda$, the torus or $B$.

\begin{lemma}[the analytic family of block operators]\label{lem:a06d:S1}
In the setting above let $B$ be real with $T_0\defeq T(B)\defeq D_B^TD_B+m_j^2+a_jQ_j(B)^TQ_j(B)\ge\gamma>0$ on $\Omega$, let $A$ be a real bond field, $s_A\defeq\sup_b|A_b|$, $w\defeq\norm qe_js_A$, $S_A$ the set of unit cubes containing an endpoint of a bond with $A_b\ne0$, and for $\tau\in\C$ let $T(\tau)$, $G(\tau)\defeq T(\tau)^{-1}$, $X(\tau)\defeq a_j-a_j^2Q_\tau G(\tau)Q_\tau^T+aL^{-2}P(\tau)$ be the operators in the background $B+\tau A$. Put $m_1\defeq5\sqrt d\gamma^{-1/2}+6a_+d\gamma^{-1}$, $c_\rho\defeq(72m_1dL)^{-1}$, $\theta_W\defeq\theta_0/4$ ($\theta_0$ the Combes--Thomas rate of group (A) at $(\gamma,1)$). For $|\tau|w\le c_\rho$: \textup{(i)} $T(\tau)$ is invertible, $G(\tau)$ analytic in $\tau$, and $\norm{T_0^{-1/2}(T(\tau)-T_0)T_0^{-1/2}}\le m_1|\tau|w\le1/72$; \textup{(ii)} $G(\tau)$ obeys the Combes--Thomas bounds of group (A) (plain, one derivative on either side, two derivatives) with rate $\theta_W$ and constant $c_{\mathrm{CT}}\defeq12\max(1,\gamma^{-1},\gamma^{-1/2})e^{\theta_W}$; \textup{(iii)} $\tilde W(\tau)\defeq X(\tau)-X(0)$ has $|\tilde W(\tau)(y,y')|\le C_W|\tau|w\sum_{\Delta\in S_A}e^{-\theta_W(\dist(y,\Delta)+\dist(\Delta,y'))}$ with $C_W=C_W(d,N,a,a_+,L,\gamma,\theta_0)$; \textup{(iv)} hence $\norm{\tilde W(\tau)}\le C_WS^2|\tau|w$ and $\sum_{y,y'}|\tilde W(\tau)(y,y')|\le C_WS^2|\tau|w|S_A|$, $S\defeq\sup_z\sum_ye^{-\theta_W\dist(y,z)}$. The same holds for every compression $\Lambda X(\tau)\Lambda$ and for $\Delta^{(j)}(\tau)$. The statement holds with $\tau A$ replaced by any complex bond field $Z$ with $w(Z)\le c_\rho$.
\end{lemma}

\begin{proof}
$U(B_b+\tau A_b)=U(B_b)U(\tau A_b)$, $\norm{U(B_b)}=1$ and $\norm{e^Z-1}\le\norm Ze^{\norm Z}$; since $|\tau|wdL\le1/72<\ln2$, $\norm{U(\tau A_b)-1}\le2\xi|\tau|w$ and $\norm{U(\tau A(\Gamma))-1}\le2\,\mathrm{len}(\Gamma)|\tau|w$. Hence $\norm{D_\tau-D_0}\le2\sqrt d|\tau|w$ (each site is the $b_+$ of at most $d$ bonds), $\norm{Q_\tau-Q_0}\le2d|\tau|w$, $\norm{Q_\tau}\le2$, $\norm{Q(\tau)-Q(0)}\le2dL|\tau|w$. (i): $T(\tau)-T_0=(D_\tau-D_0)^TD_\tau+D_0^T(D_\tau-D_0)+a_j[(Q_\tau-Q_0)^TQ_\tau+Q_0^T(Q_\tau-Q_0)]$; with $\norm{D_0T_0^{-1/2}}\le1$ and $\norm{T_0^{-1/2}}\le\gamma^{-1/2}$ the three groups contribute at most $m_1|\tau|w$, so $T(\tau)=T_0^{1/2}(1+K(\tau))T_0^{1/2}$ with $\norm K\le1/72$. (ii): conjugate by $e^{\theta\rho_f}$ as in the Combes--Thomas argument; the $\tau=0$ part contributes $\le1/8$ at $\theta\le\theta_0/4$, the $\tau$-part $V_\theta(\tau)\defeq e^{\theta\rho_f}(T(\tau)-T_0)e^{-\theta\rho_f}$ has every factor connecting points at distance $\le1$, so each norm grows by at most $e^\theta\le e$, and $\norm{e^{\theta\rho}D_\tau e^{-\theta\rho}T_0^{-1/2}}\le3$, giving $\le9m_1|\tau|w\le1/8$; so $\norm{K_\theta(\tau)}\le1/4$ and the three bounds follow with $\norm{T_\theta(\tau)^{-1}}\le\frac43\gamma^{-1}$. (iii): $\tilde W(\tau)=-a_j^2[(Q_\tau-Q_0)G_\tau Q_\tau^T+Q_0(G_\tau-G_0)Q_\tau^T+Q_0G_0(Q_\tau-Q_0)^T]+aL^{-2}[P(\tau)-P(0)]$ with $G_\tau-G_0=-G_\tau(T(\tau)-T_0)G_0$. $(Q_\tau-Q_0)^T\delta_y$ vanishes unless $B(y)\in S_A$, then is supported in $B(y)$ with norm $\le2d|\tau|w$, and $Q_\tau^T\delta_{y'}$ is supported in $B(y')$ with norm $\le2$; by (ii) the entry is $\le4dc_{\mathrm{CT}}e^{\theta_W}|\tau|we^{-\theta_W\dist(B(y),y')}$, and the third term is symmetric. $T(\tau)-T_0=\sum_{\Delta\in S_A}V_\Delta$, $V_\Delta$ collecting the bonds with $b_-\in\Delta$ and the pins of the block $\Delta$; for one $\Delta$, e.g.\ $\langle g,G_\tau(D_\tau-D_0)^T\mathbf 1_\Delta D_\tau G_0g'\rangle=\langle\mathbf 1_\Delta(D_\tau-D_0)G_\tau g,\mathbf 1_\Delta D_\tau G_0g'\rangle$, bounded by $2\sqrt d|\tau|w\norm{\mathbf 1_{\Delta^+}G_\tau g}\norm{\mathbf 1_\Delta D_\tau G_0g'}$ with $\Delta^+$ the union of the $\le3^d$ unit cubes touching $\Delta$, and the localised $L^2$ forms of (ii) bound the two factors by $3^dc_{\mathrm{CT}}e^{\theta_W}e^{-\theta_W\dist(\Delta,B(y))}\norm g$ and $c_{\mathrm{CT}}e^{-\theta_W\dist(\Delta,B(y'))}\norm{g'}$; with $g=Q_0^T\delta_y$, $g'=Q_\tau^T\delta_{y'}$ this is the $\Delta$-term. $P(\tau)-P(0)$ is block-diagonal over $L$-blocks, vanishes on $L$-blocks avoiding $S_A$, with entries $\le5dL|\tau|w\le5dLe^{2\theta_WL}|\tau|we^{-\theta_W(\dist(y,\Delta)+\dist(\Delta,y'))}$ for a $\Delta$ in the same $L$-block. Summing, $C_W\defeq a_+^2[8dc_{\mathrm{CT}}e^{2\theta_W}+4\cdot3^dc_{\mathrm{CT}}^2e^{\theta_W}(4\sqrt d+4a_+d)]+5adL^{-1}e^{2\theta_WL}$. (iv): Schur test and the sum of (iii). For a complex $Z$ the same entry bounds hold since $U$ is entire.
\end{proof}

\begin{lemma}[two kernel tools]\label{lem:a06d:G0}
\textup{(i)} Let $B$ be real with (Cur) and $Z,Z'$ complex with $w(Z),w(Z')\le c_\rho$. Then $|(X(B+Z)-X(B+Z'))(y,y')|\le C'_W\sum_\Delta\hat w_\Delta(Z-Z')e^{-\theta_W(d(y,\Delta)+d(\Delta,y'))}$, $C'_W\defeq3^de^{2\theta_W}C_W$. \textup{(ii)} Put $K_X\defeq c_X\sum_{n\in\Z^3}|n|e^{-\delta_X|n|/2}$, $S''\defeq e^2\sup_y\sum_\Delta e^{-(\theta_W/2)d(y,\Delta)}$, $\theta_2\defeq\min\{\delta_X/2,\gamma_X/(8K_X),\theta_W/2,1\}$ and $c_*\defeq\min\{c_\rho,\gamma_X/(8C'_WS''^2)\}$. For real $B$ with (Cur), complex $Z$ with $w(Z)\le c_*$, every $s\ge0$ and every $1$-Lipschitz $\rho$ on the unit lattice, $\norm{e^{\theta_2\rho}(s+X(B+Z))^{-1}e^{-\theta_2\rho}}\le2/(s+\gamma_X)$, hence $|(s+X(B+Z))^{-1}(y,y')|\le\frac2{s+\gamma_X}e^{-\theta_2|y-y'|}$.
\end{lemma}

\begin{proof}
(i) Assign each bond to the unit cube of its $b_-$, $Z-Z'=\sum_\Delta Y_\Delta$, order the cubes and set $Z_i\defeq Z'+\sum_{i'\le i}Y_{\Delta_{i'}}$; $Z_i$ agrees with $Z$ or $Z'$ bond by bond, so $w(Z_i)\le c_\rho$. Telescope $X(B+Z)-X(B+Z')=\sum_i[X(B+Z_i)-X(B+Z_{i-1})]$; the $i$-th term is $X(B+Z_{i-1}+\tau Y_{\Delta_i})-X(B+Z_{i-1})$ at $\tau=1$, bounded by \cref{lem:a06d:S1}(iii) at the complex base $B+Z_{i-1}+\tau Y_{\Delta_i}$, integrated in $\tau$; only the bonds of $Y_{\Delta_i}$ change, with amplitude $\le\hat w_{\Delta_i}(Z-Z')$, and the cubes containing their endpoints are the $\le3^d$ neighbours of $\Delta_i$, costing $e^{\theta_W}$ each. (ii) Schur test with $|e^{\theta|n|}-1|\le\theta|n|e^{\theta|n|}$: $\norm{e^{\theta\rho}X(B)e^{-\theta\rho}-X(B)}\le\theta K_X\le\gamma_X/8$; by (i) with $Z'=0$ and a Schur test, $\norm{e^{\theta\rho}(X(B+Z)-X(B))e^{-\theta\rho}}\le C'_WS''^2w(Z)\le\gamma_X/8$; $s+X(B)\ge s+\gamma_X$ is self-adjoint, so $e^{\theta\rho}(s+X(B+Z))e^{-\theta\rho}=(s+X(B))^{1/2}(1+K)(s+X(B))^{1/2}$ with $\norm K\le\frac12$. Entries: take $\rho\defeq d(\cdot,y')$.
\end{proof}
